\begin{scanpage}{0198}{180}
\noindent transmitting vibrator [compare Equation (28.03)]. The unknown coefficients $\bar A_1$ and $\bar A_2$ are found from Conditions (28.06) and equal
\[
\left.\begin{aligned}
\bar A_1&=\frac1D e^{i(k+w)z_2}\bigl[\psi_+(L)-\psi_-(L)\psi(L)e^{i(k-w)L}\bigr]e^{-ikz_1},\\
\bar A_2&=\frac1D e^{-i(k-w)z_1}\bigl[\psi_-(L)-\psi_+(L)\psi(L)e^{i(k+w)L}\bigr]e^{ikz_2}.
\end{aligned}\right\}\tag{28.17}
\]
Using equality (28.09), Expression (28.14) may be written in the more graphic form
\[
\left.\begin{aligned}
J(z)=S\bigl\{&e^{iwz}-\psi_-(z-z_1)e^{iwz_1+ik(z-z_1)}\\
&+\psi_-(L)e^{iwz_1+ikL}\bigl[\psi(z_2-z)e^{ik(z_2-z)}\\
&\qquad-\psi(L)e^{ikL}\psi(z-z_1)e^{ik(z-z_1)}\\
&\qquad+\psi^2(L)e^{2ikL}\psi(z_2-z)e^{ik(z_2-z)}-\cdots\bigr]\\
&-\psi_+(z_2-z)e^{iwz_2+ik(z_2-z)}\\
&+\psi_+(L)e^{iwz_2+ikL}\bigl[\psi(z-z_1)e^{ik(z-z_1)}\\
&\qquad-\psi(L)e^{ikL}\psi(z_2-z)e^{ik(z_2-z)}\\
&\qquad+\psi^2(L)e^{2ikL}\psi(z-z_1)e^{ik(z-z_1)}-\cdots\bigr]\bigr\}.
\end{aligned}\right\}\tag{28.18}
\]
Here besides the wave $Se^{iwz}$ and the primary edge waves, which we talked about in connection with Equation (28.14), the secondary, tertiary, etc.\ waves diverging from the ends $z=z_1$ and $z=z_2$ are explicitly written out; they correspond to the first, second, etc.\ terms in the graphs.

Passing to the limit in Equation (28.14) when $z_2\to\infty$, we find the current in the semi-infinite conductor $(z_1,\infty)$
\[
J(z)=S\cdot\bigl[e^{iwz}-\psi_-(z-z_1)e^{iwz_1+ik(z-z_1)}\bigr]\tag{28.19}
\]
and, similarly, we find the current in the semi-infinite conductor $(-\infty,z_2)$
\[
J(z)=S\cdot\bigl[e^{iwz}-\psi_+(z_2-z)e^{iwz_2+ik(z_2-z)}\bigr].\tag{28.20}
\]
\end{scanpage}
